%% The line mask, reported in two layers: the frequency coincidence, which is
%% a velocity criterion and nothing else, and the one species we do not read
%% as emission, which is an argument and is set out as one.
%% Owner: r14-mask.  Carries \label{app:linemask} (cited by 04_method.tex and
%% app_C_spatialnorm.tex) and \label{tab:masklines}.
\section{The line mask}
\label{app:linemask}

A crossing that falls on a known molecular transition has an astrophysical
explanation to hand, and the mask exists to say which crossings those are. It
reports a frequency coincidence first and an astrophysical judgement about
the coinciding species separately, so that the reproducible part can be
recomputed and the arguable part argued with.

\emph{The list, and when it was made.} The mask holds every transition of the
\FrNSpec{} species these spectral setups cover --- the three CO isotopologues,
HCN, HCO$^{+}$, CS, CN, SiO, H$_2$CO, SO, [C\,\textsc{i}] and hydrogen
recombination --- that falls inside one of the \McNIsland{} searched frequency
intervals with an upper-state energy below \McEuMax\,K. That rule is the whole
of it, and it admits \FrNTrans{} transitions (Table~\ref{tab:masklines}).
No strength cut is applied across species, since any threshold admitting the
strong lines of HCN would discard the brightest lines in these fields; within
a species the main hyperfine components are kept and the weak satellites
dropped. The list is a single Splatalogue query of CDMS and
JPL \citep{Pickett1998,CDMS2005} made on \MkListDate{}, after the disposition
criteria were fixed on \MkCritDate, and frozen with the data release, so it
is reproducible as a file rather than as a description. It is built from what
the setups cover and not from what each system is thought able to carry.

\emph{The frame is the star's own, and the sign is not the usual one.} A
crossing's frequency is carried from the topocentric value to the barycentre
using its own block's pointing and mid-time, and then to the stellar rest
frame using the star's systemic velocity. The barycentric term alone spans
\FrBarySwing\,km\,s$^{-1}$ across the survey, so a tolerance applied in the
observed frame would spend most of its budget on the Earth's motion. The
\FrEpochNEpoch{} \FrEpochStar{} epochs show
the transform working in the only way that can be checked:
\FrEpochTopoKms\,km\,s$^{-1}$ apart as observed, they agree to
\FrEpochStelKms\,km\,s$^{-1}$ in the star's frame. One row looks as though no
transform had been applied, and the reason is arithmetic: toward
\MkCancStar{} the barycentric term is \MkCancVBary\,km\,s$^{-1}$ and the
systemic velocity \MkCancVSys\,km\,s$^{-1}$, so each moves the frequency by
\MkCancBaryMHz\,MHz and the two nearly cancel, leaving \MkCancNetkHz\,kHz;
its coincidence, \MkCancDv\,km\,s$^{-1}$ from \MkCancLine, is unaffected.
Throughout, $\Delta v_\star = c\,(\nu_\star-\nu_0)/\nu_0$, so a positive
offset means the crossing lies \emph{above} the transition in frequency ---
the opposite sign to the velocity convention of a spectral line.

\begin{table}
\centering
\footnotesize
\caption{The line mask, by species: the number of transitions $N$ it holds,
the frequency range they span inside the searched intervals and their
upper-state energies. Hydrogen recombination lines carry no tabulated
$E_{\rm u}$. The transition-by-transition list is in the data release.}
\label{tab:masklines}
%% GENERATED by maskso_v414.py -- do not hand-edit.
\begin{tabular}{@{}lrr@{\,}c@{\,}rr@{}}
\hline
Species & $N$ & \multicolumn{3}{c}{$\nu_0$ (GHz)} & $E_{\rm u}$ (K) \\
\hline
H & 5 & 92.03 & -- & 231.90 & -- \\
SO & 15 & 100.03 & -- & 682.68 & 16--143 \\
H$_2$CO & 13 & 101.33 & -- & 491.97 & 10--106 \\
C$^{18}$O & 3 & 109.78 & -- & 329.33 & 5--32 \\
CN & 50 & 113.12 & -- & 680.30 & 5--114 \\
CO & 3 & 115.27 & -- & 345.80 & 6--33 \\
CS & 3 & 195.95 & -- & 342.88 & 24--66 \\
SiO & 2 & 217.10 & -- & 347.33 & 31--75 \\
$^{13}$CO & 8 & 220.40 & -- & 330.59 & 16--32 \\
HCN & 7 & 354.50 & -- & 354.51 & 43--43 \\
HCO$^{+}$ & 1 & 356.73 & -- & 356.73 & 43 \\
{[C\,\textsc{i}]} & 1 & 492.16 & -- & 492.16 & 24 \\
\hline
\end{tabular}%

\end{table}

\emph{The coincidences, and the one species we do not read as emission.} Of
the \NCross{} crossings, \MkNCoinc{} fall within
$\pm\MaskHalfKms$\,km\,s$^{-1}$ of a transition of the list in their own
star's frame; that count is the first layer and all the velocity criterion
decides. Coincidences with sulphur monoxide account for \MkNSODecide{} ---
\MkSOStars, at \MkSODvList\,km\,s$^{-1}$ from transitions of upper-state
energy \MkSOEuLo--\MkSOEuHi\,K, each inside a Band~7 window that also covers
CO($3\to2$) --- and two things argue against reading those as emission.
Sulphur monoxide has not been reported in a debris disc, where the
second-generation gas is CO, C and O; and toward \FrCpStar{} the CO($3\to2$)
line falls in the same window and is absent to \FrCpCoLimMJy\,mJy in a
\FrCpCoChanKms\,km\,s$^{-1}$ channel while the crossing itself is
\FrCpFluxMJy\,mJy, so an SO line there would have to outshine an undetected
CO($3\to2$) line in the same data. Both arguments are about the species and
neither about any one crossing, so all \MkNSODecide{} are treated alike: the
offset is printed beside each in Table~\ref{tab:ledger}, where a dagger marks
the rows this applies to, and disposes of none of them, giving
\MkAttr{} attributed against \MkUnattr{} unattributed, of
which \MkRankLeadWord{} outrank every control position in their own windows.
Reading them as attributions instead gives \MkAttrAlt{} and \MkUnattrAlt,
and \MkRankLeadAlt{} such crossing; the recurrence test excludes a persisting
carrier for all \MkNSODecide{} either way, at
\MkAttrExclLo--\MkAttrExclHi$\sigma$, so the choice moves counts and no
conclusion. The reserved hold-out holds \HoNCoinc{} coincidences,
\HoNNoted{} of them the same species and treated the same way, leaving
\HoNAttr{} attributed and \HoNUnattr{} unattributed there
(\S\ref{sec:holdoutsearch}).

\emph{The half-width, and what it costs.} A one-channel test is too tight for
circumstellar gas, whose emission is shifted by many channels: the width
follows from Keplerian motion where these molecules survive,
$\lesssim$13\,km\,s$^{-1}$ beyond 10\,au, and is
$\pm\MaskHalfKms$\,km\,s$^{-1}$. At that width the mask removes
\MkMaskAGHz\,GHz of the \CvUnionA\,GHz Class~A union, \MkMaskAPct{} per cent
of it, against \FrMaskLostGHz\,GHz of the \FrMaskUnionGHz\,GHz both classes
cover, \FrMaskPct{} per cent: the cost falls on the primary experiment,
because its windows are the line-tuned ones. Between $\pm\FrLadderLo$ and
$\pm\FrLadderHi$\,km\,s$^{-1}$ the mask labels \FrAttrWTwenty{} to
\FrAttrWHundred{} of the crossings, and the \FrNRankPass{} that outrank all
\NCtrl{} control positions of their own windows are the same two at every
width.

\emph{A bright line costs its whole window, and the windows were therefore
searched again.} Because the search records one statistic per window, a
weaker carrier in a window whose largest value is a molecular line could not
have been recorded at all (\S\ref{sec:technosearch}), so the price of line
confusion is the window and not the masked channels in it: \MkNCoinc{}
Class~A windows spanning \MkLineDomGHz\,GHz of unique frequency,
\MkLineDomPct{} per cent of the Class~A union. Most of that is recoverable,
because the per-channel statistic those maxima were taken from is retained:
removing the masked channels and taking the largest value of what is left
re-searches \MkNResearch{} of them, leaving only the \MkNWholly{} fine
windows that are narrower than the mask itself and whose \MkWhollyGHz\,GHz,
\MkWhollyPct{} per cent of the Class~A union, is lost outright. The surviving
maxima run \MkResidLo{} to \MkResidHi, so \MkNResidTrigWord{} of them
reaches the trigger: \MkSurvStar{} at \MkSurvFreq\,GHz,
\MkSurvDv\,km\,s$^{-1}$ from the nearest transition. It is what noise gives
at this threshold, since \MkNCtrlTrig{} of the \MkNCtrlPos{} control
positions of the same windows
reach it outside the mask as well, in \MkNWinCtrlTrig{} of them; and it does
not recur, the \MkSurvNRep{} other blocks covering that cell reading
$T_\star=\MkSurvTRepLo$ to \MkSurvTRepHi.

\emph{The labelled crossings are tested for recurrence too.} Of the
\MkAttr{} attributed crossings, \MkAttrTested{} have a later block covering
the predicted cell and \MkAttrRecur{} are present again, toward
\MkAttrRecurStars{} at $T_\star=\MkAttrRecurTLo$ to \MkAttrRecurTHi{} --- the
test working in the positive direction on emission that is certainly real,
and no exclusion is formed for those rows, because what it would exclude has
just been observed. For the \MkAttrExcl{} not recovered a repeat window's own
largest statistic bounds the predicted cell from above, and
Table~\ref{tab:ledger} prints the row as a bound; \MkAttrUntest{} have no
covering block at all (\MkAttrUntestNames).

\emph{Image-sideband leakage accounts for no crossing.} The mask holds
signal-sideband frequencies, and the \MkImgN{} unattributed crossings that
share a block with a CO($2\to1$) crossing in the opposite sideband would need
a first local oscillator outside the range their own block's spectral windows
permit, \MkImgLOLo--\MkImgLOHi\,GHz; their image frequencies fall at least
\MkImgDvMin\,km\,s$^{-1}$ from every masked transition.

\emph{A longer list is not a better one.} Every Splatalogue transition over
these intervals with an Einstein coefficient above $10^{-5}$\,s$^{-1}$ ---
\MaskDbTrans{} transitions of \MaskDbCarriers{} species --- would mask
\MaskDbPctUnion{} per cent of the union and move \FrFullNFlip{} crossings
across the line, \FrFullNFlipRankWord{} of them rank-leading. At these
frequencies almost anything lies within $\pm\MaskHalfKms$\,km\,s$^{-1}$ of
something in such a list, which is why the species are restricted to what
these setups cover and why the outcome is a label rather than a verdict.
